% SOURCE_STATEMENT: "For the spectral geometry of $SU_q(2)$, what are the geometric roles of Paschke duality and the ideal $K_q$? Can the Hochschild-cycle axiom be satisfied or suitably modified? Is there a canonical local-index cocycle from the spectral data? What differential calculus and metric arise, and what freedom in the Dirac operator changes that metric? What is the correct equivariant cyclic-homology framework, and how can it be computed?"
% Preparation and final self-audit: CONTRIBUTING.md prompts C and D; source/audit artifacts retained in the local physics_statements_2026_10_11 batch.
% !TeX program = lualatex
\documentclass[11pt]{article}
\usepackage[margin=25mm]{geometry}
\usepackage{fontspec}
\setmainfont{Latin Modern Roman}
\usepackage{amsmath,amssymb,mathtools}
\usepackage{unicode-math}
\setmathfont{Latin Modern Math}
\usepackage{xurl,hyperref}
\hypersetup{colorlinks=true,urlcolor=blue}
\setcounter{secnumdepth}{0}
\setlength{\parindent}{0pt}
\setlength{\parskip}{5pt}
\setlength{\emergencystretch}{3em}
\newcommand{\sref}[2]{\hyperlink{source:#1:#2}{[S#2]}}
\newcommand{\cataloguescope}[1]{\paragraph{Recorded target.}#1}
\newcommand{\R}{\mathbb R}
\newcommand{\C}{\mathbb C}
\newcommand{\N}{\mathbb N}
\newcommand{\Z}{\mathbb Z}
\newcommand{\Q}{\mathbb Q}
\newcommand{\F}{\mathbb F}
\newcommand{\eps}{\varepsilon}
\begin{document}
\section{Spectral Geometry and Cyclic Homology of $SU_q(2)$}\label{statement}

\subsection{Definitions and mathematical statement}

\paragraph{Shared notation.}$\mathbb N=\{1,2,\ldots\}$; the standard number fields have their usual meanings. All additional conventions are stated locally.


\paragraph{Definitions and conventions.} For $0<q<1$, $\mathcal A(SU_q(2))$ denotes the polynomial algebra of the compact quantum-group deformation of $SU(2)$. A spectral triple $(\mathcal A,\mathcal H,D)$ represents $\mathcal A$ on a Hilbert space with a selfadjoint operator $D$, compact resolvent, and bounded commutators $[D,a]$ on the chosen algebra. The source also uses a real structure $J$ and an equivariant quantum-group action. Its ideal $K_q$ is a compact-operator ideal generated by operators whose spectra decay exponentially; the commutant and first-order axioms hold only modulo this ideal. Hochschild cycles express an orientation axiom, while cyclic cocycles encode index pairings. The metric associated to $D$ is read from the commutator seminorm $\|[D,a]\|$.

\paragraph{Source scope and qualifications.} Keep all bundled questions: Paschke duality and K\_q; orientation; a canonical local-index cocycle; differential calculus and Dirac-operator freedom; equivariant cyclic theory. The source also asks to understand a suitable smooth algebra and calculate its cyclic homology. These are distinct requirements, not one interchangeable index computation.

\paragraph{Additional source context.} The report first asks for a satisfactory $C^\infty$ completion of the polynomial algebra and its cyclic homology. In the local-index question the desired geometric three-cocycle is distinguished from the already computed one-cocycle image.

\paragraph{Catalogue mathematical target.}
\paragraph{Catalogue wording retained for traceability.}
For the spectral geometry of $SU_q(2)$, what are the geometric roles of Paschke duality and the ideal $K_q$? Can the Hochschild-cycle axiom be satisfied or suitably modified? Is there a canonical local-index cocycle from the spectral data? What differential calculus and metric arise, and what freedom in the Dirac operator changes that metric? What is the correct equivariant cyclic-homology framework, and how can it be computed?

\paragraph{Saved literature context.} The supplied catalogue's assessment dated 2026-08-21 records: This is a bundled list of open questions in q-geometry concerning Hochschild cycles, local index data, calculi, metrics, and equivariant cyclic homology; no single general resolution was located. The assessment is retained as dated provenance; it is not a new comprehensive literature or open-status certification.
\paragraph{Scope and source qualification.}Keep all bundled questions: Paschke duality and K_q; orientation; a canonical local-index cocycle; differential calculus and Dirac-operator freedom; equivariant cyclic theory. The source also asks to understand a suitable smooth algebra and calculate its cyclic homology. These are distinct requirements, not one interchangeable index computation.
\subsection{Short English statement}
Complete the geometric, index, differential, metric, and equivariant homology interpretation of the spectral data of the deformed special unitary group.
\subsection{Sources}
\begin{itemize}
\item[\hypertarget{source:statement:1}{S1}] ULAM.AI and the UnsolvedMath Contributors, UnsolvedMath, catalogue entry. Catalogue curation under CC BY 4.0; definitions and English exposition added locally.
\url{https://huggingface.co/datasets/ulamai/UnsolvedMath}.
\item[\hypertarget{source:statement:2}{S2}] Noncommutative Geometry and Quantum Field Theory (2005), Oberwolfach Reports, DOI 10.4171/owr/2005/48, pp. 2705--2760; official report source and mathematical context.
\url{https://doi.org/10.4171/owr/2005/48}.
\item[\hypertarget{source:statement:3}{S3}] Official EMS Press PDF of the same Oberwolfach report; PDF pages 44--45, printed pages 2748--2749.
\url{https://ems.press/content/serial-article-files/46022}.
\end{itemize}
\label{end:statement}
\end{document}
